\writeSectionFrame{\presentationTitle}{Das Proxy-Muster}
\begin{frame}{Steckbrief}
	\begin{itemize}
 		\item Deutscher Name: Proxy
 		\item Alternativer deutscher Name: Stellvertreter, Surrogat
 		\item Englischer Name: Proxy
 		\item Gruppe: Strukturmuster
 	\end{itemize}
\end{frame}

\begin{frame}{Proxy}
	\begin{block}{Zweck}
 		Kontrolliere den Zugriff auf ein Objekt mithilfe eines vorgelagerten Stellvertreterobjekts.
 	\end{block}
\end{frame}

\begin{frame}{Allgemeines Klassendiagramm}
    \begin{tikzpicture}[scale=0.7, transform shape]
        \umlclass[x=6, y=6]{Subject}{}{
            +request() : void
        }
    
        \umlclass[x=6, y=1]{RealSubject}{}{
        }
    
        % Receiver
        \umlclass[x=0, y=1]{Proxy}{}{
            +action() : void
        }
    
        % Invoker
        \umlclass[, x=0, y=6]{Client}{}{
        }
    
        % Relationships
        \umlinherit{RealSubject}{Subject}
        \umlinherit{Proxy}{Subject}
        \umluniassoc{Client}{Subject}
        \umluniassoc{Proxy}{RealSubject}
    \end{tikzpicture}
\end{frame}

\begin{frame}{Mailer}
	\begin{center}
 		\includegraphics[scale=0.45]{\BILDERPFAD/proxy1.png}
 	\end{center}
\end{frame}

\note{
	Unser Programm soll automatisch Emails versenden. Das ist an sich nichts neues
	und das Programm läuft auch schon genauso. 
}

\begin{frame}{Mailer}
	\begin{center}
 		\includegraphics[scale=0.3]{\BILDERPFAD/proxy2.png}
 	\end{center}
\end{frame}

\note{
	In bestimmten Fällen soll das Versenden der Emails jedoch noch einmal separat
	protokolliert werden, bspw. wenn Rechnungen oder Mahnungen versendet werden. 
	Außerdem sollen bestimmte Emails nur zu bestimmten Zeiten verwendet werden, bspw. 
	nur während den üblichen Bürozeiten und nicht in der Nacht. Wir haben bereits
	eine Klasse Mailer, aber diese möchten wir nicht verändern, denn die neuen Funktionen
	sollen ja nicht immer eingesetzt werden. 
}

\begin{frame}{Klassendiagramm -- Mailer}
    \begin{tikzpicture}[scale=0.7, transform shape]
        \umlclass[x=6, y=6]{MailerInterface}{}{
            +sendEmail()
        }
    
        \umlclass[x=6, y=1]{Mailer}{}{
        }
    
        % Receiver
        \umlclass[x=0, y=1]{MailerProxy}{}{
        }
    
        % Invoker
        \umlclass[, x=0, y=6]{Client}{}{
        }
    
        % Relationships
        \umlinherit{Mailer}{MailerInterface}
        \umlinherit{MailerProxy}{MailerInterface}
        \umluniassoc{Client}{MailerInterface}
        \umluniassoc{MailerProxy}{Mailer}
    \end{tikzpicture}
\end{frame}

\begin{frame}[fragile]{Subjekt}
	\begin{exampleblock}{MailerInterface.java}
        \begin{lstlisting}
public interface MailerInterface {
	void sendEmail(String receiver, String sender, 
			String subject, String body);
}          
        \end{lstlisting}
 	\end{exampleblock}
\end{frame}

\begin{frame}[fragile]{Reales Subjekt}
	\begin{exampleblock}{Mailer.java}
        \begin{lstlisting}

public class Mailer implements MailerInterface {
	@Override
	public void sendeEmail(String empfaenger, String receiver, 
			String subject, String body) {
		
		System.out.println("Email wird gesendet an: " 
            + receiver + " von: " + sender);
		System.out.println("Betreff: " + subject);
		System.out.println("Inhalt: " + body);
	}
}
        \end{lstlisting}
 	\end{exampleblock}
\end{frame}

\begin{frame}[fragile]{Proxy}
	\begin{exampleblock}{TimerMailerProxy.java}
        \begin{lstlisting}
public class TimerMailerProxy implements MailerInterface {
	private static final int FIRST_SEND_HOUR = 7;
	private static final int LAST_SEND_HOUR = 18;
	
	private Mailer mailer;
	private List<MailRequest> messageQueue;
	private Random random;
	
	public TimerMailerProxy() {
		mailer = new Mailer();
		messageQueue = new ArrayList<>();
		random = new Random();
	}            
        \end{lstlisting}
 	\end{exampleblock}
\end{frame}

\begin{frame}[fragile]{Proxy}
	\begin{exampleblock}{TimerMailerProxy.java}
        \begin{lstlisting}
@Override
public void sendeEmail(String empfaenger, 
        String sender, String subject, String body) {
    int hour = random.nextInt(24);
    
    if (hour < FIRST_SEND_HOUR || hour > LAST_SEND_HOUR) {
        MailRequest mailRequest = new MailRequest(sender,
            empfaenger, body, subject);
        messageQueue.add(mailRequest);
    } else {
        mailer.sendeEmail(empfaenger, sender, subject, body);
        if (messageQueue.size() > 0) {
            for (MailRequest mailRequest: messageQueue) {
                mailer.sendeEmail(/*...*/);
            }
            messageQueue.clear();
        }
    } 
        \end{lstlisting}
 	\end{exampleblock}
\end{frame}

\begin{frame}{Klassendiagramm -- Virtuelles Dateisystem}
    \begin{tikzpicture}[scale=0.7, transform shape]
        \umlclass[x=6, y=6]{CommandExecutor}{}{
            +executeCommand(command: String)
        }
    
        \umlclass[x=6, y=1]{CommandExecutorImpl}{}{
        }
    
        % Receiver
        \umlclass[x=0, y=1]{CommandExecutorProxy}{}{
        }
    
        % Invoker
        \umlclass[, x=0, y=6]{Client}{}{
        }
    
        % Relationships
        \umlinherit{CommandExecutorImpl}{CommandExecutor}
        \umlinherit{CommandExecutorProxy}{CommandExecutor}
        \umluniassoc{Client}{CommandExecutor}
        \umluniassoc{CommandExecutorProxy}{CommandExecutorImpl}
    \end{tikzpicture}
\end{frame}

\begin{frame}[fragile]{Subjekt}
	\begin{exampleblock}{CommandExecutor.java}
        \begin{lstlisting}
public interface CommandExecutor {
	void executeCommand(String command);
}
    
        \end{lstlisting}
 	\end{exampleblock}
\end{frame}

\begin{frame}[fragile]{Reales Subjekt}
	\begin{exampleblock}{CommandExecutorImpl.java}
        \begin{lstlisting}
public class CommandExecutorImpl implements CommandExecutor {
	private VirtualFileSystem virtualFileSystem;
	
	@Override
	public void executeCommand(String command) {
		if (command.startsWith("rm")) {
			String[] parts = command.split(" ");
			String filename = parts[1];
			virtualFileSystem.rm(filename);
		} else if (command.startsWith("ls")) {
			virtualFileSystem.ls().forEach(file -> System.out.println(file));
		} else if (command.startsWith("cat")) {
			// ...;
		}
	}
}
        \end{lstlisting}
 	\end{exampleblock}
\end{frame}

\begin{frame}[fragile]{Proxy}
	\begin{exampleblock}{CommandExecutorProxy.java}
        \begin{lstlisting}
public class CommandExecutorProxy implements CommandExecutor {
	private boolean admin;
	private CommandExecutor commandExecutor;
	
	public CommandExecutorProxy(String username, String password) {
		if (username.equals("stl") && password.equals("geheim")) {
			admin = true;
		} else {
			admin = false;
		}
		
		commandExecutor = new CommandExecutorImpl();
	}     
        \end{lstlisting}
 	\end{exampleblock}
\end{frame}

\begin{frame}[fragile]{Proxy}
	\begin{exampleblock}{CommandExecutorProxy.java}
        \begin{lstlisting}
@Override
public void executeCommand(String command) {
    if (admin == true) {
        commandExecutor.executeCommand(command);
    } else if (admin == false) {
        if (command.startsWith("rm")) {
            throw new IllegalArgumentException("...");
        } else {
            commandExecutor.executeCommand(command);
        }
    }
}
        \end{lstlisting}
 	\end{exampleblock}
\end{frame}

\begin{frame}{Klassendiagramm -- Internet-Proxy}
    \begin{tikzpicture}[scale=0.7, transform shape]
        \umlclass[x=6, y=6]{Internet}{}{
            +connectTo(serverUrl: String)
        }
    
        \umlclass[x=6, y=1]{RealInternet}{}{
        }
    
        % Receiver
        \umlclass[x=0, y=1]{Proxy}{}{
        }
    
        % Invoker
        \umlclass[, x=0, y=6]{Client}{}{
        }
    
        % Relationships
        \umlinherit{RealInternet}{Internet}
        \umlinherit{Proxy}{Internet}
        \umluniassoc{Client}{Internet}
        \umluniassoc{Proxy}{RealInternet}
    \end{tikzpicture}
\end{frame}

\begin{frame}[fragile]{Subjekt}
	\begin{exampleblock}{Internet.java}
        \begin{lstlisting}
public interface Internet {
	void connectTo(String serverUrl);
}
        \end{lstlisting}
 	\end{exampleblock}
\end{frame}

\begin{frame}[fragile]{Reales Subjekt}
	\begin{exampleblock}{RealInternet.java}
        \begin{lstlisting}
public class RealInternet implements Internet {
	@Override
	public void connectTo(String serverUrl) {
		System.out.println("Ich verbinde mit " + serverUrl);
		System.out.println("******************************");
		System.out.println("******************************");
		System.out.println("******************************");
		System.out.println("Verbunden! Viel Spass");
	}
}
        \end{lstlisting}
 	\end{exampleblock}
\end{frame}

\begin{frame}[fragile]{Proxy}
	\begin{exampleblock}{InternetProxy.java}
        \begin{lstlisting}
public class InternetProxy implements Internet {
	private List<String> bannedPages;
	private Internet internet;
	
	public InternetProxy() {
		bannedPages = new ArrayList<>();
		bannedPages.add("www.freeporn.de");
		bannedPages.add("www.bomben-anleitungen.de");
		bannedPages.add("www.illegaler-waffenhandel.de");
		bannedPages.add("www.andere-gesperrte-seite.de");
		
		internet = new RealInternet();
	}
        \end{lstlisting}
 	\end{exampleblock}
\end{frame}

\begin{frame}[fragile]{Proxy}
	\begin{exampleblock}{InternetProxy.java}
        \begin{lstlisting}
@Override
public void connectTo(String serverUrl) {
    if (bannedPages.contains(serverUrl)) {
        throw new IllegalArgumentException("... " + serverUrl);
    }
    
    internet.connectTo(serverUrl);
}
        \end{lstlisting}
 	\end{exampleblock}
\end{frame}

\begin{frame}{Klassendiagramm -- Webservice}
    \begin{tikzpicture}[scale=0.7, transform shape]
        \umlclass[x=6, y=6]{WebserviceConnector}{}{
            +getItems(): List<String>
        }
    
        \umlclass[x=6, y=1]{RealWebserviceConnector}{}{
        }
    
        % Receiver
        \umlclass[x=0, y=1]{LocalWebServiceConnector}{}{
        }
    
        % Invoker
        \umlclass[, x=0, y=6]{Client}{}{
        }
    
        % Relationships
        \umlinherit{RealWebserviceConnector}{WebserviceConnector}
        \umlinherit{LocalWebServiceConnector}{WebserviceConnector}
        \umluniassoc{Client}{WebserviceConnector}
        \umluniassoc{LocalWebServiceConnector}{RealWebserviceConnector}
    \end{tikzpicture}
  \pause

  \begin{alertblock}{Bemerkung}
    Dieses Beispiel ruft einen Webservice ab, der auf einem Server im internen Stl-Netz läuft, d.h. der Proxy wird tätig, wenn Sie nicht auf den Server zugreifen können, also von außerhalb des Stl-Netzes.
  \end{alertblock}
\end{frame}

\begin{frame}[fragile]{Subjekt}
	\begin{exampleblock}{WebserviceConnector.java}
        \begin{lstlisting}
public interface WebserviceConnector {

	List<String> getItems();

}
        \end{lstlisting}
 	\end{exampleblock}
\end{frame}

\begin{frame}[fragile]{Reales Subjekt}
	\begin{exampleblock}{RealWebserviceConnector.java}
        \begin{lstlisting}
public class RealWebserviceConnector implements WebserviceConnector {
	private static final String WEBSERVICE_URL = "...";
	
	@Override
	public List<String> getItems() {
		List<String> items = new ArrayList<>();
		
		try {
            // ...

        } catch (Exception e) {
            throw new WebserviceConnectionException();
        }
		
		return items;
	}
}
        \end{lstlisting}
 	\end{exampleblock}
\end{frame}

\begin{frame}[fragile]{Proxy}
	\begin{exampleblock}{LocalWebServiceConnector.java}
        \begin{lstlisting}
public class LocalWebServiceConnector 
        implements WebserviceConnector {
	@Override
	public List<String> getItems() {
		WebserviceConnector webserviceConnector = null;
		
		try {
			webserviceConnector = new RealWebserviceConnector();
			List<String> items = webserviceConnector.getItems();
			System.out.println("Vom Webservice gelesen");
			return items;
		} catch (WebserviceConnectionException ex) {
			webserviceConnector = new LocalWebServiceConnector();
			System.out.println("Lokal gelesen");
			List<String> listWithItems = Arrays.asList(...);
			return listWithItems;
		}
	}
}
        \end{lstlisting}
 	\end{exampleblock}
\end{frame}
